% ===================================================================
%  TABLICA Nr 2 — Ciało ludzkie. — Przerwa. — Gry.
%  Meme regime que le tableau 1 : traduction de 2026, cles de
%  texto/io, renvois aux memes objets, gras sur le terme seul.
%
%  DEUX MOTS QUE LE POLONAIS SEPARE LA OU L'IDO LES CONFOND, et il
%  faut les tenir : « ramię » est le bras, « bark » l'epaule ; l'ido
%  dit « brakio » et « shultro », le francais bras et epaule, mais le
%  polonais emploie couramment « ramię » pour les deux et l'on
%  perdrait la distinction du releve. Les renvois (30) et (10) portent
%  chacun sur l'un des deux.
%
%  ET « broda » EST LA BARBE ET LE MENTON. Le mentono de l'ido est
%  « podbródek » ici, sans quoi la phrase dirait que la barbe cache la
%  barbe. Le vangobarbo est « bokobrody », le labiobarbo « wąsy ».
%
%  LES NOMS DE JEUX SONT NATIONAUX, et la note (*) du fac-simile le
%  dit elle-meme : l'ido les a forges. On prend le nom polonais quand
%  le jeu existe — kręgle, ciuciubabka, klasy —, et une description
%  quand il n'existe pas, jamais un calque.
% ===================================================================

%%K t02-apar-1 apar
\VUsaut{4.0mm}
\VUtitre{15.0pt}{60}{TABLICA Nr 2}
\VUsaut{2.0mm}
\VUtitre{11.4pt}{0}{Ciało ludzkie. --- Przerwa.}
\VUtitre{11.4pt}{0}{Gry.}

%%K t02-tit-0 sub
\VUtitre{13.2pt}{0}{Ciało ludzkie.}
\VUsaut{2.4mm}
\VUcentre{10.2pt}{0}{\textit{(Prosta lekcja nauk przyrodniczych.)}}
\VUsaut{3.0mm}

%%K t02-tit-1 sub pos=t02-01-1
\VUpk{10.2pt}{40}{Głowa.}
\VUsaut{3.0mm}

%%K t02-01-1 p
1. --- Główne części ciała ludzkiego to: \VUgras{głowa}, \VUgras{tułów} i
\VUgras{kończyny}.

%%K t02-02-1 p
2. --- Głowa zawiera dwie części: \VUgras{czaszkę}, która mieści \VUgras{mózg} i którą
okrywają \VUgras{włosy}, oraz \VUgras{twarz}.

%%K t02-03-1 p
3. --- Na naszym rysunku nie widzimy ani czaszki, ani mózgu; ale bardzo wyraźnie
widzimy ważne części twarzy.

%%K t02-04-1 p
4. --- Oto najpierw \VUgras{czoło}, które wskazuje potężny umysł, myśl zawsze czynną.
\VUgras{Skronie} są bardzo swobodne. \VUgras{Oko} jest żywe i szeroko otwarte. Gęsta
\VUgras{brew} i długie \VUgras{rzęsy} je osłaniają.

%%K t02-05-1 p
5. --- Oto dalej \VUgras{nos} i \VUgras{nozdrza}. Nos służy nam do wąchania i do
oddychania. Po obu stronach nosa są \VUgras{policzki}, dalej \VUgras{uszy}. Dolną część
twarzy tworzą \VUgras{usta} i \VUgras{podbródek}.

%%K t02-06-1 p
6. --- Nie widzimy ich zbyt dobrze, bo \VUgras{wąsy} i \VUgras{bokobrody} zakrywają je
przed nami. Ale wiemy, że usta mają dwie \VUgras{wargi}, \VUgras{szczękę} i
\VUgras{żuchwę} z ich \VUgras{zębami}, oraz \VUgras{język}. Ustami mówimy i jemy.
Językiem odczuwamy smak pokarmów.

%%K t02-07-1 p
7. --- Oko, ucho, nos i usta są, tak samo jak palce, narządami pięciu zmysłów:
\VUgras{wzroku}, \VUgras{słuchu}, \VUgras{węchu}, \VUgras{smaku}, \VUgras{dotyku}.

%%K t02-08-1 p
8. --- Głowa jest więc jedną z najciekawszych części ciała ludzkiego.

%%K t02-tit-2 sub
\VUsaut{4.4mm}
\VUpk{10.2pt}{40}{Tułów.}
\VUsaut{3.0mm}

%%K t02-09-1 p
9. --- \VUgras{Szyja} łączy głowę z tułowiem. Zawiera \VUgras{gardło} i \VUgras{kark}.

%%K t02-10-1 p
10. --- Tułów zawiera także wiele ważnych narządów. W \VUgras{piersi} są
\VUgras{płuca}, którymi oddychamy, i \VUgras{serce}, które jest ośrodkiem życia. Kiedy
serce przestaje bić, istota żywa umiera.

%%K t02-11-1 p
11. --- \VUgras{Żebra} podtrzymują pierś.

%%K t02-12-1 p
12. --- Pozostałe części tułowia to \VUgras{bok}, \VUgras{plecy}, \VUgras{barki} i
\VUgras{brzuch}. Kładziemy się na boku albo na plecach, żeby spać, nosimy ciężary na
naszych barkach.

%%K t02-13-1 p
13. --- W brzuchu są \VUgras{żołądek} i \VUgras{jelita}. Służą nam do trawienia
pokarmów.

%%K t02-tit-3 sub
\VUsaut{4.4mm}
\VUpk{10.2pt}{40}{Kończyny.}
\VUsaut{3.0mm}

%%K t02-14-1 p
14. --- Mamy cztery \VUgras{kończyny}: dwa \VUgras{ramiona} i dwie \VUgras{nogi}.
Ramionami bierzemy, trzymamy, podnosimy, zbieramy przedmioty; nogami stoimy, chodzimy i
biegamy.

%%K t02-15-1 p
15. --- \VUgras{Bark} łączy ramię z ciałem. Bardzo mocny \VUgras{mięsień}, dwugłowy,
porusza ramieniem, które \VUgras{łokieć} łączy z \VUgras{przedramieniem}.
\VUgras{Nadgarstek} tworzy staw \VUgras{dłoni}, która kończy się \VUgras{palcami} i
\VUgras{paznokciami}. Na dłoni rozróżniamy \VUgras{żyłę} (i tętnicę), w której płynie
\VUgras{krew}.

%%K t02-16-1 p
16. --- Części nogi to: \VUgras{udo}, \VUgras{kolano} i \VUgras{podkolanie},
\VUgras{łydka}, której \VUgras{mięso} jest okryte \VUgras{skórą}, \VUgras{kostka} i
\VUgras{stopa}. \VUgras{Kości} nogi są bardzo grube i bardzo mocne. Na końcu stopy są
\VUgras{palce}. Pozostałe części to \VUgras{podeszwa} i \VUgras{pięta}.

%%K t02-17-1 p
17. --- Taki jest prawie zupełny opis ciała ludzkiego.

%%K t02-17-2 p
\textit{Uwaga.} --- \textit{Za pomocą wyrażenia: « boli mnie... (głowa itd.) »
nauczyciel będzie mógł użyć całego tego słownictwa jeszcze raz, pod kątem dobrego lub
złego} \VUgras{stanu ciała}.

%%K t02-tit-4 sub
\VUsaut{5.0mm}
\VUpk{10.2pt}{40}{Gimnastyka.}
\VUsaut{2.4mm}
\VUcentre{10.2pt}{0}{\textit{(Opowiadanie jednego z uczniów.)}}
\VUsaut{3.0mm}

%%K t02-18-1 p
18. --- Żeby być giętkimi, zwinnymi i zdrowymi, musimy ćwiczyć nasze nogi i nasze
ramiona. Dlatego bawimy się i uprawiamy \VUgras{gimnastykę}.

%%K t02-19-1 p
19. --- W ciągu tygodnia te dwa ćwiczenia odbywają się na dziedzińcu liceum (zobacz
tablicę nr 1). Ale w czwartki i w niedziele idziemy na wielki \VUgras{trawnik}, gdzie
mamy miejsce, żeby biegać i się bawić.

%%K t02-20-1 p
20. --- Nasi wychowawcy i nasi rodzice czasem nam tam towarzyszą; ale ćwiczeniami
kieruje \VUgras{nauczyciel gimnastyki}\textsuperscript{(12)}.

%%K t02-21-1 p
21. --- Na trawniku, na lewo od wejścia, znajdują się \VUgras{przyrządy
gimnastyczne}\textsuperscript{(1)}, wspinamy się po \VUgras{drabinie}\textsuperscript{(2)},
robimy przewroty na \VUgras{kółkach}\textsuperscript{(3)} i na
\VUgras{trapezie}\textsuperscript{(4)}, podciągamy się na rękach aż na szczyt
\VUgras{drabinki linowej}\textsuperscript{(5)} albo \VUgras{liny z
węzłami}\textsuperscript{(6)}.

%%K t02-22-1 p
22. --- Tymczasem nasi koledzy skaczą przez \VUgras{konia
drewnianego}\textsuperscript{(7)} albo przez \VUgras{poręcze
równoległe}\textsuperscript{(11)}. Inni \VUgras{boksują}\textsuperscript{(8)} albo
\VUgras{fechtują}\textsuperscript{(9)} w maskach i \VUgras{floretami}. Wreszcie inni
\VUgras{ćwiczą z ciężarkami}\textsuperscript{(10)}.

%%K t02-tit-5 sub
\VUsaut{4.6mm}
\VUpk{10.2pt}{40}{Gry.}
\VUsaut{3.0mm}

%%K t02-23-1 p
23. --- Wokół gimnastykujących chłopcy i dziewczynki bawią się w rozmaite
\VUgras{gry}. Dziewczynki bawią się swoją \VUgras{obręczą}\textsuperscript{(14)};
skaczą przez \VUgras{skakankę}\textsuperscript{(13)} albo zapraszają swoich braci i
kuzynów na partię \VUgras{krokieta}\textsuperscript{(15)}. Cieszą się, kiedy odbiją
\VUgras{kulę} przeciwnika w chwili, gdy myśli, że łatwo przejdzie przez bramkę!

%%K t02-24-1 p
24. --- Chłopcy wolą gry bardziej siłowe. Chętnie grają w
\VUgras{kręgle}\textsuperscript{(16)}, które zręcznie przewracają
\VUgras{kulą}\textsuperscript{(17)}. Nie gardzą też zabawą
\VUgras{bąkiem}\textsuperscript{(18)} i \VUgras{bilbokietem}\textsuperscript{(19)} ani
nawet \VUgras{żabką}\textsuperscript{(20)}. Ale ich ulubione gry to
\VUgras{kulki}\textsuperscript{(21)} i \VUgras{piłka o ścianę}\textsuperscript{(22)},
bo mur \VUgras{szklarni}\textsuperscript{(26)} bardzo dobrze nadaje się do odbijania
lekkiej piłki.

%%K t02-25-1 p
25. --- Mały Jan, idący na swoich \VUgras{szczudłach}\textsuperscript{(24)}, jest
bardzo zręczny i bardzo dobrze umie utrzymać równowagę. Obok niego kilku kolegów bawi
się w \VUgras{chowanego}\textsuperscript{(25)} w tej szklarni i za \VUgras{żywopłotem}.
Inni wolą \VUgras{zgadywankę z klepnięciem}\textsuperscript{(28)} albo
\VUgras{lotkę}\textsuperscript{(29)}, która przelatuje z jednej \VUgras{rakiety} na
drugą.

%%K t02-noto-1 noto
\VUnotes{143.2mm}{%
(*) Gry dziecięce i chłopięce mają często nazwy czysto narodowe. Staraliśmy się nazwać
je w ido jak najlepiej, albo natchnąwszy się samą czynnością, która je znamionuje, albo
wybierając nazwę narodową z tego lub innego kraju, kiedy nie jest zbyt niesłuszna. Np.:
\VUgras{gra w beczkę} (po niemiecku i po francusku) to \VUgras{żabka}
\textsuperscript{(20)} (po angielsku), w której grający starają się wrzucić swoje
okrągłe krążki przez usta metalowej żaby, stojącej z otwartym pyskiem na dziurawej
skrzyni (beczce). \VUgras{Skok przez barki}\textsuperscript{(30)} różni się od
\VUgras{skoku przez plecy} tylko tym: żeby przeskoczyć przez głowę, gracze opierają ręce
na barkach towarzysza, podczas gdy w « \VUgras{skoku przez plecy} » opierają ręce tylko
na jego plecach. --- Albo też kilku uczestników jest pochylonych, a inni będą skakać
przez ich plecy opierając ręce na barku; pierwsi muszą znieść uderzenie bez ugięcia się,
drudzy muszą skoczyć bez utraty równowagi (zobacz samą tablicę).%
}

%%K t02-26-1 p
26. --- Pośrodku trawnika są dwa drzewa. U stóp jednego z nich siedmiu albo ośmiu
chłopców urządziło partię \VUgras{skoku przez barki}\textsuperscript{(30)} (*). Nie we
wszystkich krajach zna się tę grę; ale wszyscy wiedzą, czym jest \VUgras{skok przez
plecy}, w który chłopcy tym razem nie grają.

%%K t02-tit-6 sub
\VUsaut{5.0mm}
\VUpk{10.2pt}{40}{Gry na wolnym powietrzu.}
\VUsaut{3.4mm}

%%K t02-27-1 p
27. \VUgras{Ciuciubabka}\textsuperscript{(31)} też jest bardzo zabawna; dziewczynki i
chłopcy na przemian zakładają sobie \VUgras{opaskę} na oczy i chwytają niezręcznych,
którzy dają się złapać.

%%K t02-28-1 p
28. --- Pośrodku trawnika oto gra bardzo francuska, ale trochę gwałtowna: to
\VUgras{niedźwiedź}\textsuperscript{(32)}, o wiele głośniejszy niż tak spokojna gra w
\VUgras{latawca}\textsuperscript{(33)} albo nawet niż angielska gra w
\VUgras{tenisa}\textsuperscript{(34)}.

%%K t02-29-1 p
29. --- Ten ostatni sport jest radością starszych uczniów. Nasi nauczyciele sami chętnie
go uprawiają. Wymaga wielkiej zręczności i zwinności i bardzo mi się podoba.
Dziewczynkom zostawiam zabawę na \VUgras{huśtawce}\textsuperscript{(35)}.

%%K t02-30-1 p
30. --- Nie lubię innej gry angielskiej, która może stać się niebezpieczna.

%%K t02-30-1b p
Mam na myśli \VUgras{piłkę nożną}\textsuperscript{(36)}, która cieszy mojego starszego
brata. Wolę naszego starego \VUgras{berka}\textsuperscript{(38)}, równie zdrowego i
naprawdę mniej brutalnego.

%%K t02-tit-7 sub
\VUsaut{4.6mm}
\VUpk{10.2pt}{40}{Jazda na rowerze i gra w balon.}
\VUsaut{3.0mm}

%%K t02-31-1 p
31. --- Chętnie też objeżdżam trawnik na \VUgras{rowerze}\textsuperscript{(37)}.

%%K t02-32-1 p
32. --- W przyszłym roku nauczę się \VUgras{jazdy konnej}\textsuperscript{(39)}. Jakże
piękny jest koń, który wdzięcznie stąpa albo kłusuje i który galopem skacze przez
przeszkody!

%%K t02-33-1 p
33. --- Latem uczymy się \VUgras{pływania}\textsuperscript{(40)}. Z rozkoszą nurkujemy
i kąpiemy się w czystej i świeżej wodzie rzeki. Czasem na koniec puszczamy papierowy
\VUgras{balon}\textsuperscript{(41)}, który majestatycznie wznosi się w powietrze,
gdy tylko napełni się ciepłym powietrzem.

%%K t02-34-1 p
34. --- Jest jeszcze wiele innych gier, ale nie skończyłbym ich wyliczać, a z tego
streszczenia widać, że w czwartki nie nudzimy się zbytnio na naszym pięknym trawniku.

%%K t02-34-2 p
Uwaga. --- \textit{Nauczyciel łatwo uzupełni ten rozdział, opisując dokładniej każdą z
najważniejszych gier.}
